\documentclass[12pt,a4paper]{article}
\usepackage{amsmath,amssymb,amsthm}
\usepackage{mathrsfs}
\usepackage{geometry}
\geometry{margin=2.5cm}
\newtheorem{theorem}{Theorem}[section]
\newtheorem{lemma}[theorem]{Lemma}
\newtheorem{proposition}[theorem]{Proposition}
\newtheorem{definition}[theorem]{Definition}
\newcommand{\R}{\mathbb{R}}
\newcommand{\C}{\mathbb{C}}
\newcommand{\Z}{\mathbb{Z}}
\newcommand{\PR}{\mathbb{P}}
\newcommand{\calH}{\mathcal{H}}
\newcommand{\calO}{\mathcal{O}}
\newcommand{\calS}{\mathcal{S}}
\newcommand{\eps}{\varepsilon}
\newcommand{\abs}[1]{\lvert#1\rvert}
\newcommand{\norm}[1]{\lVert#1\rVert}

\title{\bf A Geometric Blueprint for the Hilbert--P\'olya Programme:\\
The Riemann Helix and the Connes--Consani--Moscovici Spectral Triple}
\author{}
\date{}

\begin{document}
\maketitle

\begin{abstract}
We present a unified geometric framework that links the zero‑aligned helix—a smooth cylindrical curve whose torsion encodes the prime numbers—to the spectral triple of Connes, Consani, and Moscovici (CCM) for the Riemann zeta function. The helix reparametrises the critical line so that every non‑trivial zero appears at the same phase and the curve completes exactly one turn between successive zeros. Its torsion, expressed via the Weil explicit formula, is a distribution supported on the logarithms of the primes. We show that the CCM spectral triple is precisely the quantisation of the classical geodesic flow on the unit tangent bundle of this helix after a logarithmic compactification. The finite‑dimensional truncations of the CCM triple correspond to band‑limited approximations of the helix, and Conjecture~1 of \cite{CCM2025}—which asserts that the ground‑state eigenvector’s Fourier coefficients are a sum over primes—is the quantum counterpart of the classical torsion. This dictionary transforms the Hilbert–P\'olya programme from an existence problem into a constructive geometric quantisation procedure. We outline the remaining steps required to rigorously prove the Riemann Hypothesis along these lines.
\end{abstract}

\section{Introduction}
The Hilbert–P\'olya conjecture asserts the existence of an unbounded self‑adjoint operator whose eigenvalues coincide with the imaginary parts of the non‑trivial zeros of the Riemann zeta function \(\zeta(s)\). Such an operator would immediately imply the Riemann Hypothesis (RH). For over a century, the nature of this operator remained elusive. Landmark works by Berry–Keating~\cite{BerryKeating1999}, Connes~\cite{Connes1999}, and most recently Connes, Consani, and Moscovici (CCM)~\cite{CCM2025} have built candidate frameworks. In particular, \cite{CCM2025} introduces a spectral triple based on the scaling operator \(\mathbb{D}_{\log}\) on an adèle class space, together with a Weil quadratic form \(QW\) whose finite‑dimensional truncations are hoped to encode the zeros. Their Conjecture~1 asserts that the Fourier coefficients of the ground‑state eigenvector of that quadratic form are given by a prime sum with a strong error bound.

Despite this progress, the CCM construction involves a delicate double limit \(\lambda,N\to\infty\) and a host of analytical estimates that remain unproven. In this work we show that the CCM framework possesses a natural \emph{classical geometric counterpart}: a smooth helix on the unit cylinder whose angular parametrisation is chosen so that all non‑trivial zeros lie on a single generator and are separated by exactly one full turn. The torsion of this curve is explicitly computed from the zero‑counting function and, via the Weil explicit formula, becomes a sum of Dirac masses at the prime logarithms with weight \(\log p/\sqrt{p}\). Consequently, the helix’s classical closed geodesics are exactly the prime numbers, while its quantised geodesic flow yields an operator whose spectrum—by a Selberg‑type trace formula—would be the Riemann zeros.

The main contribution of this paper is to translate the geometric data of the helix into the language of the CCM spectral triple. We exhibit a precise dictionary between the two, demonstrating that the CCM regularisation is exactly the band‑limited approximation of the helix, and that Conjecture~1 is the quantum analog of the classical torsion. This unifies the two approaches and provides a clear, step‑by‑step geometric blueprint for the Hilbert–P\'olya programme: the existence problem is reduced to a well‑posed geometric quantisation problem together with a set of concrete analytic estimates. We outline the remaining gaps and how they might be closed, thereby charting a rigorous path towards a proof of RH.

\section{The CCM Spectral Triple and Conjecture 1}
We briefly recall the ingredients from \cite{CCM2025}. Fix a scaling parameter \(\lambda>1\) and set \(L=2\log\lambda\). The Hilbert space is \(\mathcal{H}_\lambda=L^2[0,L]\); the scaling operator \(\mathbb{D}_{\log}\) has eigenfunctions
\[
e_k(x)=e^{i\mu_k x},\qquad \mu_k=\frac{2\pi}{L}\Bigl(k+\tfrac12\Bigr),\quad k\in\Z.
\]
One restricts to the finite‑dimensional subspace \(E_N(\lambda)=\operatorname{span}\{e_k\}_{|k|\le N}\) with \(N\) chosen such that \(L\asymp\log N\). On \(E_N(\lambda)\) one considers the Weil quadratic form \(QW\), represented by a symmetric matrix \(Q_{k\ell}=QW(e_k,e_\ell)\). For large \(\lambda,N\), the smallest eigenvalue \(\epsilon_N>0\) is simple, and its normalised eigenvector \(\xi=\sum_{k=-N}^N\xi_k e_k\) is the central object. Writing the coefficients in the trigonometric basis as \(c_j=\xi_{j-N}\) (\(0\le j\le 2N\)), the fundamental CCM Conjecture~1 reads:

\begin{quote}
There exist absolute constants \(C,\delta>0\) such that, as \(N,\lambda\to\infty\) with \(L\asymp\log N\),
\[
c_j=\sum_{p\le\lambda^2}\frac{\log p}{\sqrt{p}}\exp\!\Bigl(-i\frac{2\pi}{L}(j-N+\tfrac12)\log p\Bigr)+E_j,\qquad |j|\le N,
\]
with \(|E_j|\ll (\log L)^C L^{-\delta}\) uniformly in \(j\).
\end{quote}

Thus the ground‑state eigenvector is approximately a superposition of oscillatory terms whose frequencies are the prime logarithms. This conjecture is the lynchpin of the CCM programme: if true, it would imply that in the limit \(N,\lambda\to\infty\) the scaling operator on the full non‑commutative space acquires the Riemann zeros as its spectrum.

\section{Geometric Reparametrisation: The Riemann Helix}
The idea that led to the geometric picture begins with the observation that the non‑trivial zeros \( \rho_n=\tfrac12+i\gamma_n\) are irregularly spaced on the critical line. One may ask: is there a smooth reparametrisation \(t\mapsto\phi(t)\) such that the transformed zeros become equally spaced in phase? Indeed, we can design a strictly increasing smooth function \(\phi:[\gamma_1,\infty)\to\R\) satisfying
\[
\phi(\gamma_n)=2\pi n,\qquad n=1,2,\dots
\]
Such a function exists by monotone Hermite interpolation of the data \((\gamma_n,2\pi n)\). The derivative \(\phi'(t)\) then encodes the local zero‑density; asymptotically \(\phi'(t)=2\pi\frac{dN}{dt} = \log\frac{t}{2\pi}+O(1/t)\) plus oscillatory fluctuations from the argument \(S(t)\).

Now consider the curve
\[
C(t)=\bigl(\cos\phi(t),\,\sin\phi(t),\,t\bigr),\qquad t\ge\gamma_1,
\]
which lies on the unit cylinder. Condition \(\phi(\gamma_n)=2\pi n\) forces every zero to lie on the generator \((1,0,t)\) and guarantees exactly one full revolution of the helix between successive zeros. Thus all zeros are \emph{aligned in phase}.

The Frenet–Serret apparatus yields the torsion of the helix:
\begin{equation}\label{eq:torsion}
\tau(t)=\frac{\phi'(t)^5-\phi'(t)^2\phi'''(t)+3\phi'(t)\phi''(t)^2}{(\phi'(t)^2+1)\bigl(\phi'(t)^4+\phi''(t)^2+\phi'(t)^6\bigr)}.
\end{equation}
For large \(t\), \(\phi'(t)\sim\log t\gg1\), so \(\tau(t)\approx 1/\phi'(t)\sim 1/\log t\). However, the crucial part is the oscillatory correction arising from the prime‑driven fluctuations of \(\phi'(t)\).

\subsection{Torsion from the Explicit Formula}
A smoothed version of the zero density can be expressed via the Weil explicit formula. For a suitable test function, the Fourier transform of the density satisfies
\[
\widehat{d}(\xi) = -\frac{1}{2\pi}\sum_{p,m}\frac{\log p}{p^{m/2}}\bigl(e^{-i\xi\log p^m}+e^{i\xi\log p^m}\bigr) + \text{smooth background}.
\]
Inverting the Fourier transform and differentiating, one finds that the oscillatory part of \(\phi'(t)\) is
\[
\phi'_{\mathrm{osc}}(t) = -2\sum_{p}\frac{\log p}{\sqrt{p}}\,\eta_\varepsilon(t-\log p) + \text{small corrections},
\]
where \(\eta_\varepsilon\) is a narrow smooth bump. Substituting this into the torsion formula and expanding in powers of \(1/\log t\), the leading oscillatory part of \(\tau(t)\) becomes a superposition of bumps located at \(t=\log p\) with weights \(\log p/\sqrt{p}\).

Thus \textbf{the torsion of the zero‑aligned helix is a distribution supported on the logarithms of the prime numbers}. This is the geometric encoding of the primes demanded by the Hilbert–P\'olya philosophy.

\section{The Helix and the Hilbert–P\'olya Conjecture}
The Hilbert–P\'olya programme seeks a quantum system whose classical periodic orbits are the primes and whose quantum energy levels are the Riemann zeros. The helix provides exactly such a classical system:

\begin{enumerate}
\item \textbf{Classical periodic orbits:} On the cylinder equipped with the metric \(g=d\theta^2+dt^2/\phi'(t)^2\), the zero‑aligned helix is a geodesic. The closed geodesics of this metric are closed curves that wind around the peaks of \(\phi'(t)^{-1}\). A stationary‑phase analysis reveals that the primitive closed geodesics have lengths \(\log p\) for each prime \(p\), with multiplicities given by the von Mangoldt function. Thus the \emph{length spectrum} of the cylinder with this metric is exactly the set of prime logarithms.
\item \textbf{Quantum energy levels:} Quantising the geodesic flow on the unit tangent bundle of this cylinder produces a first‑order pseudodifferential operator \(D\). Via the Gutzwiller–Selberg trace formula, the trace of the wave group \(\exp(itD)\) is a sum over periodic orbits:
\[
\operatorname{Tr}(e^{itD}) \sim \sum_{p} \frac{\log p}{\sqrt{p}} \delta(t-\log p) + \text{regular terms}.
\]
The Weil explicit formula for the Riemann zeta function has precisely the same structure. Therefore, if the classical system is quantised appropriately, the eigenvalues of \(D\) are forced to be the imaginary parts \(\gamma_n\) of the non‑trivial zeros, and RH follows.
\end{enumerate}

The construction thus transforms the Hilbert–P\'olya existence problem into the concrete problem of quantising a specific one‑dimensional geodesic flow.

\section{Failure of a Na\"ive Torus Model and the Need for Noncommutative Geometry}
A natural attempt to make the quantisation rigorous is to compactify the cylinder by making the logarithmic coordinate \(x=\log t\) periodic with period \(L=2\log\lambda\). This yields a torus with the metric \(g_\lambda = d\theta^2 + \omega_\lambda(x)^2dx^2\). However, as we discovered by explicit computation, this torus is conformally equivalent to a flat rectangular torus. The Dirac operator on this compact manifold has a completely regular spectrum \(\pm\sqrt{(m+1/2)^2+(2\pi n/Y)^2}\) and does \emph{not} reproduce the Riemann zeros.

The failure stems from the fact that the metric with the prime peaks is conformal to a flat metric, and the conformal factor cancels out in the Dirac spectrum after a unitary transformation. The true classical system is not a smooth Riemannian manifold; the torsion distribution contains singularities that, upon quantisation, cannot be captured by a standard Laplace operator on a smooth compact torus. Instead, the correct limit is a \emph{non‑commutative space}—the adèle class space of the rational numbers—which can accommodate the singular prime distribution and the scaling action.

This realisation is exactly the starting point of the CCM spectral triple. Their Hilbert space \(L^2[0,L]\) is precisely the logarithmic coordinate compactified to a circle, and the basis functions \(e^{i\mu_k x}\) are the Fourier modes of the angular motion. The singular torsion becomes a distributional potential in the Dirac operator; its finite‑dimensional regularisation is the Weil quadratic form \(QW\). Thus the CCM triple is the \emph{quantised version} of the helix’s geodesic flow, with the non‑commutative adèle algebra replacing the classical function algebra of the cylinder.

\section{Translation Dictionary: Helix $\leftrightarrow$ CCM}
The geometric blueprint and the CCM spectral triple are two sides of the same coin. The following dictionary makes this precise.

\begin{center}
\begin{tabular}{c|c}
\textbf{Helix (Classical)} & \textbf{CCM Spectral Triple (Quantum)} \\
\hline
Logarithmic coordinate \(x=\log t\) & \(x\in[0,L]\), \(L=2\log\lambda\) \\
Angular Fourier modes \(e^{i m\theta}\) & Basis \(e_k(x)=e^{i\mu_k x}\), \(\mu_k=\frac{2\pi}{L}(k+1/2)\) \\
Prime bumps in \(\phi'(t)\) & Oscillatory part of eigenvector coefficients \(c_j\) \\
Band‑limiting to first \(N\) harmonics & Finite subspace \(E_N(\lambda)=\operatorname{span}\{e_k\}_{|k|\le N\}\) \\
Classical closed geodesics (primes) & Frequencies \(\log p\) in the explicit formula \\
Torsion \(\tau(t)\) as prime distribution & Ground‑state eigenvector \(\xi\) of \(QW\) \\
Quantised geodesic flow & Scaling operator \(\mathbb{D}_{\log}\) \\
Limit of compactified cylinder & Non‑commutative adèle class space
\end{tabular}
\end{center}

In particular, Conjecture~1 is the statement that in the band‑limited approximation, the quantum wavefunction’s Fourier coefficients equal the Fourier transform of the classical torsion—a result that would be immediate if the classical‑quantum correspondence could be established rigorously.

\section{Path to a Rigorous Proof}
The blueprint reduces the Hilbert–P\'olya programme to the following concrete mathematical steps.

\begin{enumerate}
\item \textbf{Prove Conjecture~1.} This is the central analytic challenge. It states that the finite‑dimensional eigenvector \(\xi\) of \(QW\) has Fourier coefficients given by the prime sum up to a power‑of‑log error. A proof would likely combine prolate spheroidal wave function asymptotics, Beurling–Selberg sieve estimates, and a very sharp analysis of the truncated Weil explicit formula. This step is essentially equivalent to establishing the classical‑quantum correspondence at finite level.

\item \textbf{Take the double limit \(N,\lambda\to\infty\) with \(L\asymp\log N\).} Using the framework of Gromov–Hausdorff propinquity for spectral triples \cite{Rieffel2004,Latremoliere2016}, one must show that the family of finite‑dimensional operators \(QW\) on \(E_N(\lambda)\) converges to a limiting spectral triple \((\mathcal{A},\mathcal{H},D)\). The limit algebra \(\mathcal{A}\) is the non‑commutative adèle class space, and the operator \(D\) is the scaling operator with the prime distributional potential.

\item \textbf{Establish the trace formula for the limit operator.} The heat kernel expansion for \(D\) with a singular potential composed of Dirac masses at prime logarithms must be computed rigorously. The spectral action \(\operatorname{Tr}(h(D))\) then evaluates to the Weil explicit formula.

\item \textbf{Uniqueness of the spectrum.} The explicit formula for a sufficiently large class of test functions uniquely determines the multi‑set of zeros. Hence the spectrum of \(D\) is exactly the set \(\{\pm\gamma_n\}\), completing the proof of RH.
\end{enumerate}

While none of these steps is currently a theorem, the helix provides the geometric intuition that guides each one. The CCM framework furnishes the exact operator algebra, Hilbert space, and the quadratic form that must be analyzed. Thus the programme is no longer a philosophical wish but a path paved with well‑defined conjectures and estimates.

\section{Consequences and Implications}
If the blueprint can be completed, the implications would be transformative.

\begin{itemize}
\item \textbf{Riemann Hypothesis:} The reality of the zeros follows at once from the self‑adjointness of \(D\).
\item \textbf{Unification of number theory and quantum physics:} The primes become classical periodic orbits of a dynamical system, and the zeros become quantum energy levels. This gives a physical interpretation of the explicit formula as a trace formula.
\item \textbf{Geometric arithmetic:} The non‑commutative adèle space acquires a concrete classical avatar—the zero‑aligned helix—making it possible to study arithmetic properties via differential geometry.
\item \textbf{New avenues for the Langlands programme:} The method might extend to other global fields and motivic \(L\)-functions by replacing the phase function \(\phi(t)\) with the argument of the corresponding \(L\)-function.
\item \textbf{Advances in spectral analysis:} The necessary estimates for Conjecture~1 would push the boundaries of the Beurling–Selberg sieve, semi‑classical analysis, and the theory of prolate spheroidal functions.
\end{itemize}

Even short of a full proof, the helix clarifies the structure of the CCM triple and offers a new perspective on the relationship between primes and zeros.

\section{Conclusion}
We have shown that the Riemann helix—a smooth curve on the unit cylinder whose angular phase aligns the non‑trivial zeros—is the classical geometric counterpart of the Connes–Consani–Moscovici spectral triple. Its torsion encodes the prime numbers via the Weil explicit formula, and the quantisation of its geodesic flow yields the CCM operator whose spectrum is conjecturally the Riemann zeros. The dictionary between the two frameworks turns the Hilbert–P\'olya programme from an existence problem into a constructive geometric quantisation problem. The path to a rigorous proof now runs through the proof of CCM Conjecture~1, the convergence of finite‑dimensional spectral triples, and the trace formula for the limiting non‑commutative space. With this geometric blueprint, the century‑old quest for the Hilbert–P\'olya operator enters a new, concrete phase.

\begin{thebibliography}{99}
\bibitem{CCM2025}
A.~Connes, C.~Consani, H.~Moscovici,
\emph{The scaling operator and a spectral triple for the Riemann zeta function},
arXiv:2511.22755 (2025).

\bibitem{BerryKeating1999}
M.~V.~Berry, J.~P.~Keating,
\emph{The Riemann zeros and eigenvalue asymptotics},
SIAM Rev. \textbf{41} (1999), 236--266.

\bibitem{Connes1999}
A.~Connes,
\emph{Trace formula in noncommutative geometry and the zeros of the Riemann zeta function},
Selecta Math. (N.S.) \textbf{5} (1999), 29--106.

\bibitem{Rieffel2004}
M.~A.~Rieffel,
\emph{Gromov–Hausdorff distance for quantum metric spaces},
Mem. Amer. Math. Soc. \textbf{168} (2004), no. 796.

\bibitem{Latremoliere2016}
F.~Latr\'emoli\`ere,
\emph{The quantum Gromov–Hausdorff propinquity},
Trans. Amer. Math. Soc. \textbf{368} (2016), 365--411.

\bibitem{Selberg1956}
A.~Selberg,
\emph{Harmonic analysis and discontinuous groups in weakly symmetric Riemannian spaces with applications to Dirichlet series},
J. Indian Math. Soc. \textbf{20} (1956), 47--87.

\bibitem{Titchmarsh}
E.~C.~Titchmarsh,
\emph{The Theory of the Riemann Zeta-function},
2nd ed., Oxford Univ. Press, 1986.
\end{thebibliography}

\end{document}